\begin{abstract}
Reward models guide RLHF training, yet aggregate evaluation metrics conceal systematic failure patterns. We introduce a $2\times2$ mode decomposition that classifies preference battles by human vote entropy (high/low) crossed with reward model variance (high/low), revealing four distinct alignment modes. Analyzing 57,477 Chatbot Arena battles, we find that \textbf{23.7\%} fall into Mode~3 (Misaligned-Confident): samples where human voters disagree while the reward model expresses high confidence. This overconfident misalignment is not a fringe phenomenon---approximately one in four preference battles exhibits this failure pattern. We test two mechanistic hypotheses: (1) that Mode~3 arises from semantic divergence between responses, and (2) that Mode~3 concentrates in subjective prompt types. Both are unsupported: Mode~3 pairs show \emph{higher} semantic similarity than aligned pairs (Cohen's $d = -0.049$), and Mode~3 rates are nearly identical across subjective and objective tasks (ratio $= 1.07$). The phenomenon exists; its mechanism remains unknown. Our contributions are: (a) the first empirical quantification of overconfident RM misalignment in large-scale preference data, (b) the mode decomposition framework as an alignment diagnostic, and (c) negative results ruling out candidate mechanisms. These findings demonstrate that aggregate metrics are insufficient for reward model evaluation and motivate mode-aware benchmarking.
\end{abstract}
